% BHU PYQ -- Home Science: Apparel Design (2025-26 NEP)
\documentclass[11pt,a4paper]{article}
\usepackage[a4paper,top=2.5cm,bottom=2.5cm,left=2.8cm,right=2.8cm,headheight=14pt]{geometry}
\usepackage{amsmath,amssymb}
\usepackage[T1]{fontenc}
\usepackage[utf8]{inputenc}
\usepackage{lmodern,microtype,fancyhdr,enumitem,hyperref,lastpage,parskip}

\hypersetup{colorlinks=true,urlcolor=black,linkcolor=black,pdftitle={HSCMJ32: Apparel Design -- BHU NEP 2025-26}}

\pagestyle{fancy}\fancyhf{}
\renewcommand{\headrulewidth}{0.4pt}\renewcommand{\footrulewidth}{0.4pt}
\fancyhead[L]{\small \textit{Banaras Hindu University}}
\fancyhead[R]{\small \textit{HSCMJ32: Apparel Design (2025-26)}}
\fancyfoot[C]{\small Page \thepage\ of \pageref{LastPage}}

\newlist{parts}{enumerate}{2}
\setlist[parts,1]{label=(\alph*),leftmargin=*,itemsep=4pt,topsep=3pt}
\setlist[parts,2]{label=(\roman*),leftmargin=*,itemsep=2pt,topsep=2pt}

\newcommand{\pts}[1]{\hfill \textit{[#1]}}

\begin{document}

\begin{center}
  {\large\bfseries BANARAS HINDU UNIVERSITY}\\[2pt]
  {\normalsize B. A. / B. Sc. (Hons.) Semester III Examination, 2025-26 (NEP)}\\[6pt]
  \rule{0.8\linewidth}{0.5pt}\\[6pt]
  {\large\bfseries HOME SCIENCE}\\[4pt]
  {\large\bfseries Paper Code: HSCMJ32 : Apparel Design}\\[4pt]
  {\normalsize Time: 2:30 Hours \hfill Full Marks: 70}\\[2pt]
  {\small REF NO. 2025-26/D-III/131}
\end{center}
\rule{\linewidth}{0.4pt}
\smallskip

\noindent \textit{Instructions: Answer all Sections. The figures in the right-hand margin indicate marks. Write your Roll No.\ at the top immediately on receipt of this question paper.}

\bigskip

\section*{SECTION -- A}
\noindent \textit{Note: Answer any \textbf{three} of the following questions (each within 250 words).}\pts{3 $\times$ 10 = 30}

\begin{enumerate}[label=\textbf{\arabic*.},leftmargin=*,itemsep=10pt]

  \item Explain the components of garment design (silhouette, texture and colour). Also discuss their applications in designing for different body figures.\pts{10}

  \item Explain the detailed procedure and precautions to be considered while taking body measurements for various garments.\pts{10}

  \item What are the different techniques of pattern making? Explain any one method with their applications and limitations.\pts{10}

  \item What do you understand by structural design details used in garment construction? Elaborate any two in detail.\pts{10}

  \item Explain the essential equipments required for measuring, marking and cutting during garment construction.\pts{10}

\end{enumerate}

\bigskip

\section*{SECTION -- B}
\noindent \textit{Note: Answer any \textbf{four} of the following questions (each within 100 words).}\pts{4 $\times$ 5 = 20}

\begin{enumerate}[label=\textbf{\arabic*.},leftmargin=*,start=6,itemsep=8pt]

  \item What are the major sources of inspiration for dress designing?\pts{5}

  \item Explain the importance of interlining, underlining and lining.\pts{5}

  \item Write about the techniques of achieving rhythm in dress.\pts{5}

  \item Write about any two common fitting problems and their remedies.\pts{5}

  \item Discuss about any five basic embroidery stitches used in applied designing.\pts{5}

\end{enumerate}

\bigskip

\section*{SECTION -- C}
\noindent \textit{Note: Write short notes about the following within 50 words each:}\pts{10 $\times$ 2 = 20}

\begin{enumerate}[label=\textbf{11.}(\alph*),leftmargin=*,itemsep=6pt]

  \item Emphasis in apparel designing\pts{2}
  \item Pattern envelope\pts{2}
  \item Slash and spread method\pts{2}
  \item Slot seam\pts{2}
  \item French seam\pts{2}
  \item Couching\pts{2}
  \item Patch work\pts{2}
  \item Hemming\pts{2}
  \item Facing\pts{2}
  \item Fasteners\pts{2}

\end{enumerate}

\end{document}
